5.1.2 Effect of Popularity Confounding

Figure~\ref{fig:rho_sweep} isolates the mechanism behind Theorem~1 and Corollary~1 in a controlled setting where the popularity-confounding strength $\rho$ is swept explicitly while the ground-truth preference ranking is held fixed. As $\rho$ increases from 0 to 4, Vanilla's representation uniformity loss $\mathcal{L}_{\mathrm{uni}}$ systematically moves toward zero (Fig.~\ref{fig:rho_sweep}a), indicating that user-context representations become increasingly concentrated as the shared, time-varying popularity signal is absorbed into them. This representation collapse is accompanied by a marked decline in ground-truth preference-ranking Recall@10 (Fig.~\ref{fig:rho_sweep}b): the model's utility-only ranking increasingly reflects popularity rather than the true user preference it was trained to recover. Ours, which routes the popularity signal through a separate Hawkes-intensity branch and enforces the centering constraint on the utility term, remains nearly invariant in both representation uniformity and Recall@10 across the full range of $\rho$. Together, the two panels provide controlled empirical evidence for the mechanism described in Theorem~1 and Corollary~1: when temporal popularity is not modeled separately, a shared popularity component can be absorbed into the utility representations and distort their geometry, and this distortion is consistent with the corresponding degradation in ranking quality. Ours mitigates this effect by explicitly separating the popularity component from the centered utility term. All Vanilla-vs-Ours differences shown here are statistically significant (paired $t$-test, $p<3.60\times10^{-49}$ for $\mathcal{L}_{\mathrm{uni}}$ and $p<6.62\times10^{-11}$ for Recall@10 across all $\rho$; full tests in Appendix~\ref{app:rho_sweep_stats}).
